\chapter{\textit{stable}-Unabhängigkeit}
\label{ch:stable}

\section{Aussagenlogische Kodierung der \textit{stable}-Semantik}
\label{sec:st-pl}

In diesem Abschnitt wird die \textit{stable}-Semantik in einer
aussagenlogischen Form kodiert. Ziel ist es, die Menge der
\textit{stable}-Labellings eines AFs durch eine aussagenlogische Formel zu
beschreiben. Diese Kodierung bildet anschließend die Grundlage für die
QBF-Formulierung der \textit{stable}-Unabhängigkeit.

Wir verwenden die folgende Charakterisierung der \textit{stable}-Semantik.
Sei \(F=(A,R)\) ein AF und sei \(S\subseteq A\). Für
\[
    S^+ := \{b\in A \mid \exists a\in S : (a,b)\in R\}
\]
ist \(S\) eine \textit{stable} Extension genau dann, wenn
\[
    S^+ = A\setminus S.
\]

Da in \textit{stable}-Labellings kein \(\mathit{undecided}\)-Label
vorkommt, genügt es, nur die \(\mathit{in}\)-gelabelten Argumente
aussagenlogisch zu kodieren. Ein Argument, dessen \(in\)-Variable falsch
ist, wird entsprechend als \(\mathit{out}\) interpretiert.

\begin{Definition}
    \label{def:in-vars}
    Sei \(F=(A,R)\) ein AF. Für jedes Argument \(a\in A\) sei \(i_a\)
    eine aussagenlogische Variable mit der Bedeutung
    \[
        i_a \text{ ist wahr genau dann, wenn } \lambda(a)=\mathit{in}.
    \]
    Für die \(k\)-te Kopie der Labelling-Variablen schreiben wir
    \[
        I^k := \{i_a^k \mid a\in A\}.
    \]
    Eine Belegung dieser Variablen kodiert ein Labelling \(\lambda_k\),
    indem gilt:
    \[
        \lambda_k(a)=\mathit{in}
        \quad\Longleftrightarrow\quad
        i_a^k \text{ ist wahr}.
    \]
    Ist \(i_a^k\) falsch, so wird \(a\) als \(\mathit{out}\) gelabelt.
\end{Definition}

\begin{Definition}
    \label{def:st-pl}
    Sei \(F=(A,R)\) ein AF und sei \(I = \{i_a \mid a\in A\}.\)

    Die aussagenlogische Kodierung der
    \textit{stable}-Semantik ist definiert durch
    \[
        \mathrm{st}_F(I)
        :=
        \bigwedge_{a\in A}
        (
            i_a
            \leftrightarrow
            \bigwedge_{(b,a)\in R} \neg i_b
        ).
    \]
    Es gilt
    \[
    \bigwedge \{\} = \top.
    \]
\end{Definition}

Die Formel besagt: Ein Argument \(a\) ist genau dann \(\mathit{in}\),
wenn keiner seiner Angreifer \(\mathit{in}\) ist. Da alle nicht
\(\mathit{in}\)-gelabelten Argumente in der \textit{stable}-Semantik als
\(\mathit{out}\) interpretiert werden, ist damit das gesamte Labelling
bestimmt.

\begin{Lemma}
    \label{lem:st-pl-correct}
    Sei \(F=(A,R)\) ein AF und sei \(\nu\) eine Belegung der Variablen
    \(I=\{i_a\mid a\in A\}\). Mit
    \[
        S_\nu := \{a\in A \mid \nu(i_a)=\top\}.
    \]
    Dann gilt:
    \[
        \nu \models \mathrm{st}_F(I)
        \quad\Longleftrightarrow\quad
        S_\nu \text{ ist eine \textit{stable} Extension von } F.
    \]
\end{Lemma}

\begin{proof}
    Sei zunächst \(\nu\models \mathrm{st}_F(I)\). Dann gilt für jedes
    \(a\in A\):
    \[
        \nu(i_a)=\top
        \quad\Longleftrightarrow\quad
        \nu(i_b)=\bot
        \text{ für alle } b \text{ mit } (b,a)\in R.
    \]
    Nach Definition von \(S_\nu\) ist dies äquivalent zu
    \[
        a\in S_\nu
        \quad\Longleftrightarrow\quad
        \text{kein Angreifer von }a\text{ liegt in }S_\nu.
    \]
    Damit ist \(S_\nu\) konfliktfrei. Außerdem gilt für jedes
    \(a\notin S_\nu\), dass \(a\) von einem Argument aus \(S_\nu\)
    angegriffen wird. Also gilt
    \[
        S_\nu^+ = A\setminus S_\nu.
    \]
    Folglich ist \(S_\nu\) eine \textit{stable} Extension von \(F\).

    Sei umgekehrt \(S_\nu\) eine \textit{stable} Extension von \(F\).
    Dann gilt \(S_\nu^+=A\setminus S_\nu\). Also ist für jedes
    \(a\in A\) genau dann \(a\in S_\nu\), wenn kein Angreifer von \(a\)
    in \(S_\nu\) liegt. Mit der Definition von \(S_\nu\) ergibt sich
    somit für jedes \(a\in A\):
    \[
        \nu(i_a)=\top
        \quad\Longleftrightarrow\quad
        \nu(i_b)=\bot
        \text{ für alle } b \text{ mit } (b,a)\in R.
    \]
    Daher erfüllt \(\nu\) jede Konjunktion von \(\mathrm{st}_F(I)\), also
    \[
        \nu\models \mathrm{st}_F(I).
    \]
\end{proof}

Die Menge \(S_\nu\) induziert eindeutig das Labelling \(\lambda_\nu\) mit
\(\lambda_\nu(a)=\mathit{in}\) für \(a\in S_\nu\) und
\(\lambda_\nu(a)=\mathit{out}\) für \(a\notin S_\nu\).
Daher kodiert \(\nu\) genau dann ein \textit{stable}-Labelling, wenn
\(\nu\models \mathrm{st}_F(I)\).

\section{QBF-Kodierung der \textit{stable}-Unabhängigkeit}
Wir verwenden die folgende Definition bedingter Unabhängigkeit für AFs nach \citeauthor{rienstra:2020:independence-af}\cite{rienstra:2020:independence-af}.

\begin{Definition}
    \label{def:BU:AFs}
    Sei $F = (A,R)$ ein AF, $\sigma$ eine Semantik und $X,Y,Z \subseteq A$ paarweise disjunkte Mengen. $X$ ist $\sigma$-unabhängig von $Y$
    gegeben $Z$ gdw. für alle $\lambda_1, \lambda_2 \in \Lambda_{\sigma}(F)$, wenn $\lambda_1|_Z=\lambda_2|_Z$, dann existiert ein
    $\lambda_3 \in \Lambda_{\sigma}(F)$ sodass $\lambda_3|_X=\lambda_1|_X$, $\lambda_3|_Y=\lambda_2|_Y$ und $\lambda_3|_Z=\lambda_1|_Z=\lambda_2|_Z$.
    Wir schreiben $X \bot_{\sigma} Y\mid Z$.
\end{Definition}

\begin{Lemma}
    \label{lem:eq-qbf}
    Sei \(F=(A,R)\) ein AF und seien
    \(\lambda_j,\lambda_k \in \Lambda_{\mathrm{st}}(F)\)
    \textit{stable}-Labellings von \(F\). Sei
    \[
        I^l=\{i_a^l \mid a\in A\}
    \]
    die zugehörige Variablenmenge des \(l\)-ten Labellings.
    Für eine Teilmenge \(S\subseteq A\) definieren wir
    \[
        \mathrm{Eq}^{j,k}_S
        :=
        \bigwedge_{a\in S}
        \left(
            i_a^j \leftrightarrow i_a^k
        \right).
    \]
    Unter Belegungen von \(I^j\cup I^k\), die \(\lambda_j\) und
    \(\lambda_k\) kodieren, gilt:
    \[
        \mathrm{Eq}^{j,k}_S \text{ ist wahr}
        \quad \Longleftrightarrow \quad
        \lambda_j|_S = \lambda_k|_S.
    \]
\end{Lemma}

\begin{proof}
    Für \(a\in S\) gilt unter den kodierenden Belegungen
    \(i_a^j\leftrightarrow i_a^k\) genau dann, wenn
    \(\lambda_j(a)=\mathit{in}\) genau dann gilt, wenn
    \(\lambda_k(a)=\mathit{in}\). Da \(\lambda_j\) und \(\lambda_k\)
    stable-Labellings sind, ist jedes nicht-\(\mathit{in}\)-Argument
    \(\mathit{out}\). Daher stimmen die Wahrheitswerte der \(i\)-Variablen
    auf \(S\) genau dann überein, wenn die Labellings auf \(S\)
    übereinstimmen.
\end{proof}

Auf Grundlage der Formeln \(st_F\) und \(\mathrm{Eq}^{j,k}_S\) lässt sich die \textit{stable}-Unabhängigkeit als \textit{QBF} ausdrücken.

\begin{Proposition}
    \label{prop:st-ind-qbf}
    Sei \(F=(A,R)\) ein AF und seien
    \(X,Y,Z\subseteq A\) paarweise disjunkte Mengen. Für
    \(k\in\{1,2,3\}\) sei
    \[
        I^k=\{i_a^k \mid a\in A\}
    \]
    eine Kopie der aussagenlogischen Labelling-Variablen.

    Dann wird die \textit{stable}-Unabhängigkeit von \(X\) und \(Y\)
    gegeben \(Z\), also \(X \bot_{\mathrm{st}} Y \mid Z\), durch die
    folgende geschlossene QBF kodiert:
    \begin{multline*}
    \Phi^{\mathrm{st}}_F(X,Y \mid Z)
    :=
    \forall I^1, I^2
    \exists I^3
    \left(
        \left(
            \mathrm{st}_F(I^1)
            \land
            \mathrm{st}_F(I^2)
            \land
            \mathrm{Eq}^{1,2}_Z
        \right)
        \right. \\
        \left.
        \rightarrow
        \left(
            \mathrm{st}_F(I^3)
            \land
            \mathrm{Eq}^{1,3}_X
            \land
            \mathrm{Eq}^{2,3}_Y
            \land
            \mathrm{Eq}^{1,3}_Z
        \right)
    \right).
    \end{multline*}
    Es gilt:
    \[
        \Phi^{\mathrm{st}}_F(X,Y \mid Z)\text{ ist wahr}
        \quad\Longleftrightarrow\quad
        X \bot_{\mathrm{st}} Y \mid Z.
    \]
\end{Proposition}

\begin{proof}
    Sei \(\Phi^{\mathrm{st}}_F(X,Y\mid Z)\) wahr. Seien
    \(\lambda_1,\lambda_2\in\Lambda_{\mathrm{st}}(F)\) beliebige
    \textit{stable}-Labellings mit
    \[
        \lambda_1|_Z=\lambda_2|_Z.
    \]
    Die zugehörigen Belegungen von
    \(I^1\) und
    \(I^2\) erfüllen dann
    \[
        \mathrm{st}_F(I^1)
        \land
        \mathrm{st}_F(I^2)
        \land
        \mathrm{Eq}^{1,2}_Z.
    \]
    Also ist das Antezedens der Implikation in
    \(\Phi^{\mathrm{st}}_F(X,Y\mid Z)\) wahr. Da die QBF wahr ist,
    existiert eine Belegung von
    \(I^3\), sodass
    \[
        \mathrm{st}_F(I^3)
        \land
        \mathrm{Eq}^{1,3}_X
        \land
        \mathrm{Eq}^{2,3}_Y
        \land
        \mathrm{Eq}^{1,3}_Z
    \]
    gilt. Nach Lemma~\ref{lem:st-pl-correct} kodiert diese Belegung ein
    \textit{stable}-Labelling
    \(\lambda_3\in\Lambda_{\mathrm{st}}(F)\). Aus den
    Gleichheitsformeln folgt
    \[
        \lambda_3|_X=\lambda_1|_X,\qquad
        \lambda_3|_Y=\lambda_2|_Y,\qquad
        \lambda_3|_Z=\lambda_1|_Z.
    \]
    Damit gilt \(X\bot_{\mathrm{st}}Y\mid Z\).

    Sei umgekehrt \(X\bot_{\mathrm{st}}Y\mid Z.\) Wir betrachten eine beliebige Belegung der universell quantifizierten Variablen
    \(I^1\) und \(I^2\). Ist das Antezedens
    \[
        \mathrm{st}_F(I^1)
        \land
        \mathrm{st}_F(I^2)
        \land
        \mathrm{Eq}^{1,2}_Z
    \]
    falsch, so ist die Implikation wahr.

    Sei also das Antezedens wahr. Dann kodieren die Belegungen von
    \(I^1\) und \(I^2\) nach Lemma~\ref{lem:st-pl-correct} zwei
    \textit{stable}-Labellings 
    \[
        \lambda_1,\lambda_2\in\Lambda_{\mathrm{st}}(F)
    \]
   Aus \(\mathrm{Eq}^{1,2}_Z\) folgt mit Lemma~\ref{lem:eq-qbf}, dass
    \[
        \lambda_1|_Z=\lambda_2|_Z.
    \]
    Wegen \(X\bot_{\mathrm{st}}Y\mid Z\) existiert ein
    \textit{stable}-Labelling
    \(\lambda_3\in\Lambda_{\mathrm{st}}(F)\), sodass
    \[
        \lambda_3|_X=\lambda_1|_X,\qquad
        \lambda_3|_Y=\lambda_2|_Y,\qquad
        \lambda_3|_Z=\lambda_1|_Z=\lambda_2|_Z.
    \]
    Wähle nun die Belegung von \(I^3\), die \(\lambda_3\) kodiert.
    Dann gilt
    \[
        \mathrm{st}_F(I^3)
        \land
        \mathrm{Eq}^{1,3}_X
        \land
        \mathrm{Eq}^{2,3}_Y
        \land
        \mathrm{Eq}^{1,3}_Z.
    \]
    Somit ist die Konsequens wahr. Da dies für jede Belegung der
    universell quantifizierten Variablen gilt, ist
    \(\Phi^{\mathrm{st}}_F(X,Y\mid Z)\) wahr.
\end{proof}

\section{Komplexität}
\begin{Corollary}
    \label{cor:ind-st-in-pi2p}
    Das Entscheidungsproblem \(\textit{Ind}_{\mathrm{st}}\) liegt in
    \(\Pi_2^\text{P}\).
\end{Corollary}

\begin{proof}
    Sei eine Instanz von \(\textit{Ind}_{\mathrm{st}}\) gegeben durch ein
    AF \(F=(A,R)\) und paarweise disjunkte Mengen \(X,Y,Z\subseteq A\).
    Nach Proposition~\ref{prop:st-ind-qbf} gilt
    \[
        X\bot_{\mathrm{st}}Y\mid Z
        \quad\Longleftrightarrow\quad
        \Phi^{\mathrm{st}}_F(X,Y\mid Z)\text{ ist wahr}.
    \]
    Damit existiert eine Reduktion von \(\textit{Ind}_{\mathrm{st}}\)
    auf das Wahrheitsproblem der QBF
    \(\Phi^{\mathrm{st}}_F(X,Y\mid Z)\).

    Die Formel \(\Phi^{\mathrm{st}}_F(X,Y\mid Z)\) hat die Quantorenstruktur \(\forall U\exists V\,\psi\)
    wobei \(U\) und \(V\) Mengen von aussagenlogischen Variablen sind und \(\psi\) eine quantorenfreie aussagenlogische Formel ist.

    Die Konstruktion von \(\Phi^{\mathrm{st}}_F(X,Y\mid Z)\) ist
    polynomial in der Größe von \(F\). Insgesamt werden \(3|A|\) Labelling-Variablen eingeführt. Die Formel
    \(\mathrm{st}_F(I^k)\) hat eine Größe, die polynomial in \(|A| + |R| \) liegt. Außerdem ist \(\mathrm{Eq}^{i,j}_S\) linear in \(|S|\).

    Das Wahrheitsproblem für QBFs der Form \(\forall U\exists V\,\psi\)
    liegt in \(\Pi_2^\text{P}\). Da jede Instanz von
    \(\textit{Ind}_{\mathrm{st}}\) in Polynomialzeit auf eine solche QBF
    reduziert werden kann, folgt
    \[
        \textit{Ind}_{\mathrm{st}} \in \Pi_2^\text{P}.
    \]
\end{proof}